\chapter{Reference Card}

\section{Loading the Disk}

\begin{lstlisting}
DIR
LOAD "COMMON"
LIST 60000-60280
LOAD "TINY-1"
RUN
\end{lstlisting}

\section{Framework Variables}

\begin{center}
\begin{tabularx}{0.88\textwidth}{>{\ttfamily}lX}
\toprule
ZT\$ & Game title.\\
ZA\$ & Author or series line.\\
ZS\$ & Subtitle.\\
ZP & Prompt mode: 0 for Enter or Space, 1 for Space.\\
\bottomrule
\end{tabularx}
\end{center}

\section{Variable Rule}

Only the first two characters of a NovaBASIC variable name are significant.
Use one- or two-character names in printed listings. Avoid \texttt{Z*} in game
code because the common framework owns that prefix.

\section{Line Map}

\begin{center}
\begin{tabularx}{0.88\textwidth}{>{\ttfamily}lX}
\toprule
10 & Program setup.\\
20-9999 & Game logic.\\
60000 & Common title framework begins.\\
60280 & Common title framework returns.\\
\bottomrule
\end{tabularx}
\end{center}
